% GENERATED FILE. Edit canonical lessons, then run npm run generate:books.
% canonical-source-hash: fnv1a64:9a34c7b0720e81bb
% canonical-lessons: GE-C80-loeffel, GE-C80-flasche, GE-C80-heft, GE-C80-stift, GE-C80-papier

\chapter{Spoon, Bottle, Exercise book, Pen, Paper}
\label{ch:ge-spoon-bottle-exercise-book-pen-paper}
\hlchaptermodality{80}

\begin{chapteropening}
\textbf{By the end of this chapter:} \emph{I can say a spoon, a bottle, an exercise book, a pen, a pencil and paper.}

Complete the last lesson of chapter 80: I can say a spoon, a bottle, an exercise book, a pen, a pencil and paper.
\end{chapteropening}

\section[der Löffel]{der Löffel — a spoon}

\label{lesson:GE-C80-loeffel}



\begin{quote}
Before the new one: say the German for a fork, then the German for a knife.
\end{quote}

\subsection*{You'll want to know: der Löffel}
\textbf{der Löffel} — \textquotedblleft{}a spoon\textquotedblright{}.

\begin{cousinweb}
\textbf{Löffel} $\leftarrow$ Old High German \emph{leffil} $\leftarrow$ Proto-Germanic \emph{*lapilaz}, \textquotedblleft{}spoon\textquotedblright{} — inherited; Dutch \emph{lepel} is the same word. (EtymDB 2.1) It is related to English \emph{lap}, \textquotedblleft{}to lick up\textquotedblright{}.
\end{cousinweb}

\begin{grammarlens}[title={What you've built}]
One more word for this chapter.
\end{grammarlens}

\subsection*{Your turn}
\begin{itemize}
\raggedright
  \item \emph{Say it:} \emph{der Löffel}
  \item \emph{Say it:} \emph{der Löffel}, once more
  \item \emph{Say it:} say \emph{der Löffel}
  \item \emph{Recall:} say the German for a fork, then the German for a knife, then say \emph{der Löffel} again
\end{itemize}

\begin{tcolorbox}[breakable,colback=teal!4,colframe=teal!35!black,title={Before you move on}]
What does der Löffel mean? (\textquotedblleft{}A spoon\textquotedblright{}.) Say it once more.
\end{tcolorbox}
\section[die Flasche]{die Flasche — a bottle}

\label{lesson:GE-C80-flasche}



\begin{quote}
Before the new one: say the German for a knife, then the German for a spoon.
\end{quote}

\subsection*{You'll want to know: die Flasche}
\textbf{die Flasche} — \textquotedblleft{}a bottle\textquotedblright{}.

\begin{cousinweb}
\textbf{Flasche} $\leftarrow$ Old High German \emph{flasca} $\leftarrow$ Proto-Germanic \emph{*flaskōn}, \textquotedblleft{}a wicker-covered bottle\textquotedblright{}, from a verb for \textquotedblleft{}to plait\textquotedblright{} — inherited. (EtymDB 2.1)
\end{cousinweb}

\begin{grammarlens}[title={What you've built}]
One more word for this chapter.
\end{grammarlens}

\subsection*{Your turn}
\begin{itemize}
\raggedright
  \item \emph{Say it:} \emph{die Flasche}
  \item \emph{Say it:} \emph{die Flasche}, once more
  \item \emph{Say it:} say \emph{die Flasche}
  \item \emph{Recall:} say the German for a knife, then the German for a spoon, then say \emph{die Flasche} again
\end{itemize}

\begin{tcolorbox}[breakable,colback=teal!4,colframe=teal!35!black,title={Before you move on}]
What does die Flasche mean? (\textquotedblleft{}A bottle\textquotedblright{}.) Say it once more.
\end{tcolorbox}
\section[das Heft]{das Heft — an exercise book}

\label{lesson:GE-C80-heft}



\begin{quote}
Before the new one: say the German for a spoon, then the German for a bottle.
\end{quote}

\subsection*{You'll want to know: das Heft}
\textbf{das Heft} — \textquotedblleft{}an exercise book\textquotedblright{}.

\begin{cousinweb}
\textbf{Heft} comes from a verb meaning \textquotedblleft{}to fasten\textquotedblright{} $\leftarrow$ Old High German \emph{heften} $\leftarrow$ Proto-Germanic \emph{*haftijaną}, \textquotedblleft{}to bind\textquotedblright{} (EtymDB 2.1): a booklet is sheets fastened together.
\end{cousinweb}

\begin{grammarlens}[title={What you've built}]
One more word for this chapter.
\end{grammarlens}

\subsection*{Your turn}
\begin{itemize}
\raggedright
  \item \emph{Say it:} \emph{das Heft}
  \item \emph{Say it:} \emph{das Heft}, once more
  \item \emph{Say it:} say \emph{das Heft}
  \item \emph{Recall:} say the German for a spoon, then the German for a bottle, then say \emph{das Heft} again
\end{itemize}

\begin{tcolorbox}[breakable,colback=teal!4,colframe=teal!35!black,title={Before you move on}]
What does das Heft mean? (\textquotedblleft{}An exercise book\textquotedblright{}.) Say it once more.
\end{tcolorbox}
\section[der Stift]{der Stift — a pen, a pencil}

\label{lesson:GE-C80-stift}



\begin{quote}
Before the new one: say the German for a bottle, then the German for an exercise book.
\end{quote}

\subsection*{You'll want to know: der Stift}
\textbf{der Stift} — \textquotedblleft{}a pen, a pencil\textquotedblright{}.

\begin{cousinweb}
\textbf{Stift} $\leftarrow$ Middle High German \emph{stift}, \textquotedblleft{}pin, spike\textquotedblright{} (EtymDB 2.1) — inherited, and probably related to the word for \textquotedblleft{}stiff\textquotedblright{}. A pencil is a pointed rod.
\end{cousinweb}

\begin{grammarlens}[title={What you've built}]
One more word for this chapter.
\end{grammarlens}

\subsection*{Your turn}
\begin{itemize}
\raggedright
  \item \emph{Say it:} \emph{der Stift}
  \item \emph{Say it:} \emph{der Stift}, once more
  \item \emph{Say it:} say \emph{der Stift}
  \item \emph{Recall:} say the German for a bottle, then the German for an exercise book, then say \emph{der Stift} again
\end{itemize}

\begin{tcolorbox}[breakable,colback=teal!4,colframe=teal!35!black,title={Before you move on}]
What does der Stift mean? (\textquotedblleft{}A pen, a pencil\textquotedblright{}.) Say it once more.
\end{tcolorbox}
\section[das Papier]{das Papier — paper}

\label{lesson:GE-C80-papier}



\begin{quote}
Before the new one: say the German for an exercise book, then the German for a pen.
\end{quote}

\subsection*{You'll want to know: das Papier}
\textbf{das Papier} — \textquotedblleft{}paper\textquotedblright{}.

\begin{cousinweb}
\textbf{Papier} $\leftarrow$ Middle High German \emph{papier} $\leftarrow$ Latin \emph{papyrus} $\leftarrow$ Greek \emph{pápyros} (EtymDB 2.1) — borrowed, like English \emph{paper}.
\end{cousinweb}

\begin{grammarlens}[title={What you've built}]
One more word for this chapter.
\end{grammarlens}

\subsection*{Your turn}
\begin{itemize}
\raggedright
  \item \emph{Say it:} \emph{das Papier}
  \item \emph{Say it:} \emph{das Papier}, once more
  \item \emph{Say it:} say \emph{das Papier}
  \item \emph{Recall:} say the German for an exercise book, then the German for a pen, then say \emph{das Papier} again
\end{itemize}

\begin{tcolorbox}[breakable,colback=teal!4,colframe=teal!35!black,title={Before you move on}]
What does das Papier mean? (\textquotedblleft{}Paper\textquotedblright{}.) Say it once more.
\end{tcolorbox}
